% =========================================================
% Monotonicity of Supremum and Infimum
% =========================================================

\subsection{Order Preservation of Supremum and Infimum}
\label{sec:bound-monotone}

% ---------------------------------------------------------
% Toolkit
% ---------------------------------------------------------
\begin{tcolorbox}[colback=gray!6, colframe=gray!40, arc=2pt,
  left=6pt, right=6pt, top=4pt, bottom=4pt,
  title={\small\textbf{Monotonicity — Quick Reference}},
  fonttitle=\small\bfseries]
\renewcommand{\arraystretch}{1.4}
\begin{tabular}{@{}p{0.28\textwidth} p{0.44\textwidth} p{0.20\textwidth}@{}}
\textbf{Result} & \textbf{Statement} & \textbf{Proof} \\
\midrule
Monotonicity of $\sup$
  & $A \subseteq B \Rightarrow \sup A \le \sup B$
  & \hyperref[prf:monotonicity-of-supremum]{Proof $\to$} \\
Monotonicity of $\inf$
  & $A \subseteq B \Rightarrow \inf B \le \inf A$
  & \hyperref[prf:monotonicity-of-infimum]{Proof $\to$} \\
Functional version
  & $A \subseteq B \Rightarrow \sup_A f \le \sup_B f$
  & immediate corollary \\
Diameter monotonicity
  & $A \subseteq B \Rightarrow \operatorname{diam}(A) \le \operatorname{diam}(B)$
  & corollary \\
\end{tabular}
\end{tcolorbox}

% ---------------------------------------------------------
\subsubsection{The Core Theorem}
% ---------------------------------------------------------

\begin{remark*}[Geometric reading]
The monotonicity result says: a smaller set cannot have a larger supremum.
This is geometrically immediate.
$\sup B$ is an upper bound for $B$.
Every element of $A$ is also in $B$, so $\sup B$ bounds $A$ as well.
Since $\sup A$ is the \emph{least} upper bound of $A$, it cannot exceed any
upper bound of $A$, and in particular $\sup A \le \sup B$.

The infimum direction is the mirror image: $\inf B$ is a lower bound for $B$,
hence for any $A \subseteq B$, and since $\inf A$ is the greatest lower bound
it satisfies $\inf A \ge \inf B$.

The figure below makes both directions visible at once.
\end{remark*}

\begin{figure}[h]
\centering
\begin{tikzpicture}[>=Stealth, thick, font=\small]
  \draw[->] (-0.5,0) -- (11.5,0) node[right] {$\mathbb{R}$};

  % Set B (larger, red)
  \draw[red!70!black, line width=5pt, line cap=round] (1.0,0) -- (9.5,0);
  \node[red!70!black, above] at (5.25,0.12) {$B$};
  \draw[red!70!black, dashed, thin] (1.0,-0.35)--(1.0,0.75)
    node[above]{\scriptsize $\inf B$};
  \draw[red!70!black, dashed, thin] (9.5,-0.35)--(9.5,0.75)
    node[above]{\scriptsize $\sup B$};

  % Set A (smaller, blue), drawn slightly below so both are visible
  \draw[blue!70!black, line width=5pt, line cap=round] (3.0,-0.18) -- (6.5,-0.18);
  \node[blue!70!black, below] at (4.75,-0.22) {$A \subseteq B$};
  \draw[blue!70!black, dashed, thin] (3.0,-0.55)--(3.0,0.35)
    node[above]{\scriptsize $\inf A$};
  \draw[blue!70!black, dashed, thin] (6.5,-0.55)--(6.5,0.35)
    node[above]{\scriptsize $\sup A$};

  % Inequality annotations
  \draw[<->, gray!60, thin] (1.0,-0.85) -- (3.0,-0.85)
    node[midway, below, font=\scriptsize, gray!80] {$\inf B \le \inf A$};
  \draw[<->, gray!60, thin] (6.5,-0.85) -- (9.5,-0.85)
    node[midway, below, font=\scriptsize, gray!80] {$\sup A \le \sup B$};
\end{tikzpicture}
\caption{$A \subseteq B$: the endpoints of $A$ lie inside $B$.
The infimum of $A$ is pushed right of $\inf B$ and the supremum of $A$
is pushed left of $\sup B$, so $\inf B \le \inf A \le \sup A \le \sup B$.}
\label{fig:monotonicity-sets}
\end{figure}

\begin{theorem}[\texorpdfstring{\hyperref[prf:monotonicity-of-supremum]{Monotonicity of Supremum and Infimum}}{Monotonicity of Supremum and Infimum}]
\label{thm:monotonicity}
Let $A,B \subseteq \mathbb{R}$ be nonempty.

\begin{enumerate}
\item \textbf{Supremum.}
If $A \subseteq B$ and both sets are bounded above, then
\[
\sup A \le \sup B.
\]

\item \textbf{Infimum.}
If $A \subseteq B$ and both sets are bounded below, then
\[
\inf B \le \inf A.
\]
\end{enumerate}
\end{theorem}

\begin{remark*}[Fully quantified form]
$\forall A, B \subseteq \mathbb{R}$ nonempty, bounded above:
$A \subseteq B \Rightarrow \sup A \le \sup B$.
$\forall A, B \subseteq \mathbb{R}$ nonempty, bounded below:
$A \subseteq B \Rightarrow \inf B \le \inf A$.
\end{remark*}

\begin{remark*}[Proof shape]
Both parts are one-line arguments.
For (1): $\sup B$ is an upper bound of $B$, hence of $A$ (since $A \subseteq B$);
$\sup A$ is the least upper bound of $A$, so $\sup A \le \sup B$.
For (2): the identical argument with all inequalities reversed.
No $\varepsilon$-argument is needed — the definition of supremum as
\emph{least} upper bound does all the work directly.
\end{remark*}

% ---------------------------------------------------------
\subsubsection{Functional Version}
% ---------------------------------------------------------

\begin{theorem}[\texorpdfstring{\hyperref[prf:monotonicity-of-supremum]{Monotonicity for Functions}}{Monotonicity for Functions}]
\label{thm:monotonicity-functions}
Let $f : X \to \mathbb{R}$ and let $A,B \subseteq X$ with $A \subseteq B$.

\begin{enumerate}
\item If $f$ is bounded above on $B$, then
\[
\sup_{x \in A} f(x)
\le
\sup_{x \in B} f(x).
\]

\item If $f$ is bounded below on $B$, then
\[
\inf_{x \in A} f(x)
\ge
\inf_{x \in B} f(x).
\]
\end{enumerate}
\end{theorem}

\begin{remark*}[Why this is just the set version]
Apply Theorem~\ref{thm:monotonicity} to the image sets
$f(A) = \{f(x) : x \in A\}$ and $f(B) = \{f(x) : x \in B\}$.
Since $A \subseteq B$, every value $f(a)$ for $a \in A$ is also
a value $f(b)$ for some $b \in B$, so $f(A) \subseteq f(B)$.
Monotonicity of supremum then gives $\sup f(A) \le \sup f(B)$,
which is the displayed inequality written in functional notation.
\end{remark*}

% ---------------------------------------------------------
\subsubsection{Consequences}
% ---------------------------------------------------------

\begin{corollary}[\texorpdfstring{\hyperref[prf:monotonicity-of-supremum]{Diameter Monotonicity}}{Diameter Monotonicity}]
\label{cor:diam-monotone}
Let $(X,d)$ be a metric space and let $A \subseteq B \subseteq X$.
Then
\[
\operatorname{diam}(A) \le \operatorname{diam}(B).
\]
\end{corollary}

\begin{remark*}[Why]
$\operatorname{diam}(A) = \sup\{d(x,y) : x,y \in A\}$.
The set of pairwise distances $\{d(x,y) : x,y \in A\}$ is a subset of
$\{d(x,y) : x,y \in B\}$ because $A \subseteq B$.
Monotonicity of supremum gives the result.
\end{remark*}

\begin{corollary}[Norm Bounds]
\label{cor:norm-bounds}
Let $A,B \subseteq \mathbb{R}^n$ with $A \subseteq B$. Then
\[
\sup_{x\in A} \|x\|
\le
\sup_{x\in B} \|x\|.
\]
\end{corollary}

\begin{corollary}[Integral Bounds]
\label{cor:integral-bounds}
Let $f : X \to \mathbb{R}$ be integrable and let $A \subseteq B$.
Then
\[
\inf_{x\in B} f(x)
\le
\int_A f
\le
\sup_{x\in B} f(x).
\]
\end{corollary}

\begin{remark*}[Connections forward]
Monotonicity of supremum and infimum appears in:
\begin{itemize}
  \item \emph{Bound arithmetic} (\S\ref{sec:bound-arithmetic}):
        the additivity and difference identities all rely on identifying
        a candidate bound and showing it is the least (or greatest) one;
        monotonicity is the tool that rules out smaller upper bounds.
  \item \emph{Darboux integration}: upper and lower sums satisfy
        $L(f,P) \le L(f,Q)$ whenever $P \subseteq Q$ (refinement), a direct
        application of infimum monotonicity to the set of lower sums.
  \item \emph{Convergence}: if $\{a_n\}_{n \ge N}$ is a tail of a bounded
        sequence then $\sup\{a_n : n \ge N\}$ is non-increasing in $N$
        (the sets are nested), which is the basis for $\limsup$.
  \item \emph{Diameter estimates} in metric spaces: the diameter corollary
        above is used when showing that nested closed sets with diameters
        tending to zero have a unique common point.
\end{itemize}
\end{remark*}
